\documentclass{article}

\usepackage{corl_2026} % Use this for the initial submission.
% \usepackage[final]{corl_2026} % Uncomment for the camera-ready ``final'' version.
% \usepackage[preprint]{corl_2026} % Uncomment for pre-prints (e.g., arxiv); This is like ``final'', but will remove the CORL footnote.

\title{FEMR: Front-End Motion Refinement for Robust Motion Tracking}

% The \author macro works with any number of authors. There are two
% commands used to separate the names and addresses of multiple
% authors: \And and \AND.
%
% Using \And between authors leaves it to LaTeX to determine where to
% break the lines. Using \AND forces a line break at that point. So,
% if LaTeX puts 3 of 4 authors names on the first line, and the last
% on the second line, try using \AND instead of \And before the third
% author name.

% NOTE: authors will be visible only in the camera-ready and preprint versions (i.e., when using the option 'final' or 'preprint'). 
% 	For the initial submission the authors will be anonymized.

\author{
  Jane E.~Doe\\
  Department of Electrical Engineering and Computer Sciences\\
  University of California Berkeley 
  United States\\
  \texttt{janedoe@berkeley.edu} \\
  %% examples of more authors
  %% \And
  %% Coauthor \\
  %% Affiliation \\
  %% Address \\
  %% \texttt{email} \\
  %% \AND
  %% Coauthor \\
  %% Affiliation \\
  %% Address \\
  %% \texttt{email} \\
  %% \And
  %% Coauthor \\
  %% Affiliation \\
  %% Address \\
  %% \texttt{email} \\
  %% \And
  %% Coauthor \\
  %% Affiliation \\
  %% Address \\
  %% \texttt{email} \\
}


\begin{document}
\maketitle

%===============================================================================

\begin{abstract}
    % 具身智能的终极目标之一, 是赋予人形机器人类似人类的 "即看即做" (See-and-do) 能力: 即直接在物理世界中零样本（Zero-shot）执行互联网或真实世界中的未知视频中的动作，而无需任何离线动力学解算

	% 但当我们试图直接部署从野外视频（In-the-wild videos）中提取的动作序列时，由于相机抖动与单目深度歧义，视觉工具链提取出的运动学序列不可避免地带有浮空与高频跃迁等视觉伪影

	% 任何预训练的底层控制策略都拥有一个隐形的鲁棒性预算. 然而, 如果游参考动作充满伪影，Tracker就会被动耗尽预算去强行维持平衡, 导致其在真实的 Sim-to-Real 部署中变得极其脆弱

	% 为了解决这个问题, 我们在缺陷参考轨迹和冻结的Tracker之间插入了一个轻量级的在线修正模块FrontRES. FrontRES在执行期会输出Delta SE对参考帧进行即时修正.

	% Moreover, 为了防止盲目追求“误差清零”而输出超出动力学承受极限的过度修正，FrontRES 通过“可执行奖励（Executable Reward）”精准学习了底层追踪器的性能边界，确保输出的修正量能抵消噪声的同时不反向破坏Tracker原有的物理执行能力

	% 实验展示了FrontRES能够提升GMT跟踪带噪参考帧的鲁棒性

	Enabling humanoid robots with 'see-and-do' capability - zero-shot execution of motions from in-the-wild videos - is a central goal in embodied intelligence.However, kinematic sequences extracted from these videos inevitably carry artifacts: floating, jitter, and contact inconsistency. These artifacts drain the tracker's \textit{robustness budget} - its finite capacity to absorb sensing errors and real-world perturbations. When this budget is spent on artifacts, the robot becomes fragile under real-world disturbance. To address this, we propose FEMR (Front-End Motion Refiner), a lightweight module placed between flawed reference and frozen tracker. FEMR predicts per-frame $\Delta SE(3)$ corrections at runtime. Crucially, FEMR is trained with an executable reward that respects the tracker's physical limits, rather than blindly minimizing kinematic error. Experiments demonstrate that FEMR substantially improves tracking robustness against noisy references, enabling stable zero-shot execution.
\end{abstract}

% Two or three meaningful keywords should be added here
\keywords{Robots, Learning} 

%===============================================================================

\section{Introduction}

    % 1、宏大愿景与数据范式 (The Paradigm)
	% 2、提出矛盾：“伪影”的致命性 (The Artifact Crisis)
	% 3、提出理论：“鲁棒性预算” (The Robustness Budget)
	% 4、提出需求: (The Requirements)
	% 5、提出方案：FEMR 的优雅架构 (Introducing FEMR)
	% 6、贡献列表 (Summary of Contributions)
	
	%%%
	% 第 1 段：通用运动追踪器改变了问题的瓶颈
	% 1、提出长期愿景：人形机器人不应只执行手工设计的技能，而应能复用开放世界中的人类动作。
	% 2、保留“得益于现代通用运动追踪器的飞速发展”这一表述，低层控制器可以跟踪多样的参考动作
	% 3、指出瓶颈转移：问题不再只是“控制器能不能完成动作”，而变成“有无多样且可靠的参考动作”
	% 4、引出野外视频：互联网视频包含极限动作和交互动作，比传统 MoCap 更丰富
	%%%

	% 赋予人形机器人复现开放世界人类动作的能力，是具身智能的重要目标之一
	
	% 得益于现代通用运动追踪器的发展，人形机器人已经能够通过统一的底层控制器跟踪多样的参考动作

	% 因此，限制机器人动作能力的瓶颈正在发生变化

	% 问题不再只是控制器能否完成给定动作，而是能否获得足够多样且可靠的参考动作

	% 传统 MoCap 数据具有较好的物理一致性，但采集成本高，动作覆盖范围有限

	% 相比之下，互联网视频包含更丰富的极限动作和交互动作，是更具扩展性的参考动作来源

	Enabling humanoid robots with open-world motion capabilities is a central goal in embodied intelligence. Benefiting from recent advances in general motion tracking~\cite{}, humanoid robots can now track diverse reference motions through a unified low-level controller~\cite{}, without per-skill retraining.
	
	As a result, the bottleneck is shifting. The key question is no longer only whether a controller can execute a given reference motion stably - it is also whether we can obtain reference motions that are both diverse and reliable enough to begin with. Traditional motion capture provides physically consistent data with clean calibration, but it is expensive to collect and limited in motion coverage. In contrast, Internet videos span a far richer space of extreme motions and interactive behaviors, making them a more scalable source of reference motions.
	
	%%%
	% 第 2 段：野外视频恢复出的参考动作并不是控制器可直接执行的参考轨迹
	% 1、指出核心困难：从单目或互联网视频中恢复出的三维人体运动，本质上只是视觉估计结果
	% 2、伪影类型：相机运动、深度歧义、人体重建和机器人重定向会带来根节点浮空/穿地、滑步、抖动等问题
	% 3、这些错误不是普通的视觉噪声，因为它们直接改变了接触、支撑和质心运动等控制器依赖的物理条件
	% 4、若强迫 tracker 跟踪这样的类人但接触和支撑并不自洽的轨迹；Tracker会在扰动下更容易摔倒
	%%%

	% 然而，从野外视频中恢复出的三维人体动作并不是为闭环物理控制准备的参考轨迹

	% 单目估计、人体重建和机器人重定向都会引入参考帧伪影
	
	% 这些伪影包括根节点浮空或穿地、脚步滑移、朝向抖动和接触时序不一致
	
	% 它们在视觉上可能并不明显
	
	% 但是，它们会直接改变接触、支撑和质心运动等物理条件
	
	% 当跟踪器被迫执行这些参考动作时，机器人会更容易失去稳定性

	However, 3D human motions recovered from in-the-wild videos are not ready-to-use references for closed-loop physical control. They are visual estimates, not physically validated trajectories. Camera shake, monocular depth ambiguity, human pose reconstruction~\cite{}, and robot retargeting~\cite{} all introduce reference-frame artifacts: root floating or penetration, foot sliding, orientation errors, and contact timing mismatch. Such artifacts can be subtle in visual space - a pose may look plausible frame by frame. In physical space, however, these artifacts directly corrupt the signals that controllers depend on: contact states, support relationships, and body balance. When a tracker is forced to follow flawed references, the robot becomes substantially more vulnerable to disturbances.
	
	%%%
	% 第 3 段：提出“鲁棒性预算”视角，说明伪影为什么致命
	% 1、我们用一个简单诊断实验揭示这种参考帧误差的后果
	% 2、当根节点参考噪声增加时，受到推力后 ZMP margin 会下降
	% 3、增加扰动强度会导致整段动作的 end-to-end success rate 下降
	% 4、说明参考帧伪影会消耗 tracker 的鲁棒性预算
	%%%

	% 我们将这一问题称为 Kinematic-to-Dynamic Gap
	
	% 它描述了运动学合理性和动力学可执行性之间的失配
	
	% 一个参考动作可以在姿态空间中看起来自然，却仍然难以被物理机器人稳定执行
	
	% 这种失配会消耗跟踪器的鲁棒性预算
	
	% 如 Fig.~\ref{fig:robustness-budget} 所示，随着参考帧噪声增大，机器人在外部推力后的稳定裕度持续下降
	
	% 整段动作的 end-to-end success rate 也随之降低
	
	% 这说明参考帧伪影本身会削弱跟踪器抵抗扰动的能力

	We characterize this mismatch as a \textit{kinematic-to-dynamic gap}: a reference motion can appear natural in pose space yet still be difficult for a physical robot to execute stably. This gap is not merely a qualitative intuition - it directly consumes what we call the \textit{robustness budget} of the tracker. Every low-level controller has a finite capacity to absorb sensing errors, modeling noise, and external perturbations. When reference-frame artifacts already occupy part of that budget, less remains for handling real-world disturbances. A diagnostic experiment (Fig.~\ref{fig:robustness-budget}) confirms this: as root reference noise increases, the robot's post-push stability margin drops monotonically, and the end-to-end success rate of the full motion sequence declines accordingly.
	
	%%%
	% 第 4 段：提出解决方案应当满足的要求
	% 1、合适的方法需要满足三个条件：
	% 2、放在 tracker 前，而不是改 tracker；
	% 3、在线轻量，不依赖离线动力学解算；
	% 4、tracker-aware，目标是提升可执行性，而不是让动作看起来更像人。
	%%%

	% 一个合适的解决方案应当满足三个条件

	% 首先，它应当放在跟踪器之前，而不是替换或重训跟踪器

	% 底层跟踪器已经学习到复杂的动力学控制结构

	% 让它同时承担参考修复和运动控制，会混合两个不同的问题

	% 其次，该方法应当在线且轻量

	% 它不能依赖离线动力学优化，也不能为每段视频重新求解控制轨迹

	% 然后，它应当针对参考动作中的空间伪影
	
	% 根节点高度、身体朝向和水平位移误差可以在单帧层面被修正
	
	% 相比之下，接触时序错位和动作节奏错误通常需要重新对齐整段动作，这并不是本文的目标

	% 最后，该方法应当是 tracker-aware

	% 它的目标不是让参考动作在视觉上更像人，而是让参考动作更容易被冻结跟踪器执行

	A suitable solution should satisfy four requirements. First, it should be placed before the tracker rather than replace or retrain it. Modern trackers already encode complex dynamic control capabilities~\cite{}; asking the same module to simultaneously repair upstream errors and handle downstream dynamics conflates two distinct problems. Second, it should be online and lightweight, requiring no offline trajectory optimization or per-video dynamics solving. Third, it should target spatial artifacts - errors in root height, body orientation, and horizontal displacement - that can be corrected at the per-frame level, rather than contact timing misalignment which demands sequence-level temporal realignment. Fourth, it should be \textit{tracker-aware}: the objective is not to make the reference look more human-like, but to make it easier for the frozen tracker to execute.
	
	%%%
	% 第 5 段：介绍 FEMR
	% 第 1 句：提出方法：FrontRES 是插入缺陷参考流和冻结 GMT 之间的前端残差修正器
	% 第 2 句：说明执行形式：在执行期，FrontRES 为每个参考帧预测 Δ SE(3) 修正和置信度
	% 第 3 句：强调核心设计：FEMR 不盲目把运动学误差清零，而是提高冻结 tracker 的可执行度
	% 第 4 句：说明训练信号：为此，FEMR 使用 executable reward 提升Tracker的可执行性
	%%%

	% 为了解决这一问题，我们提出 FEMR，一个放置在缺陷参考动作和冻结运动跟踪器之间的前端运动修正器

	% FEMR 不学习新的底层控制器

	% 它只在参考动作进入跟踪器之前进行轻量修正

	% 具体来说，FEMR 在执行时为每个参考帧预测任务空间中的 $\Delta SE(3)$ 修正量

	% 同时，FEMR 预测位置和姿态修正的置信度

	% 这种设计使 FEMR 能够控制何时修正、修正哪些维度以及修正幅度

	% 当参考动作已经可执行时，FEMR 应保持输入不变

	% 当参考帧伪影破坏跟踪稳定性时，FEMR 应输出有效修正

	We propose FEMR (Front-End Motion Refiner), a lightweight module inserted between corrupted reference motions and a frozen tracker. FEMR is not a new motion tracker. Instead, it predicts per-frame task-space $\Delta SE(3)$ corrections - residual adjustments to root translation and orientation - gated by separate confidence scores for position and orientation. These confidence scores give FEMR a critical property: it defaults to restraint. When the reference motion is already executable, FEMR keeps the input unchanged. When artifacts threaten stability, FEMR outputs corrections just large enough to restore executability, without overshooting into dynamically infeasible commands. To learn this behavior, FEMR is trained with a tracker-aware executable reward that measures whether a correction actually improves physical execution, rather than merely reducing kinematic error.
	
	%%%
	% 第 6 段：贡献列表
	% 第 1 句：贡献 1：将参考帧伪影表述为通用运动追踪中的 Kinematic-to-Dynamic Gap
	% 第 2 句：贡献 2：提出一个轻量在线前端修正器，在保持 GMT 冻结的同时避免离线动力学解算
	% 第 3 句：贡献 3：提出 tracker-aware executable reward 和鲁棒性预算诊断用于训练
	% 第 4 句：贡献 4：进行了实验验证FEMR的有效性
	%%%

	% 本文的贡献如下

	% 我们将视频参考动作中的伪影表述为通用运动追踪中的 Kinematic-to-Dynamic Gap，并从鲁棒性预算的角度分析其影响

	% 我们提出 FEMR，一个在线轻量的前端运动修正器。该方法在保持底层跟踪器冻结的同时，提高缺陷参考动作的可执行性

	% 我们提出 tracker-aware 的可执行性优化项设计，使 FEMR 学习何时保持不变以及何时输出有效修正
	
	% 我们通过扰动恢复实验验证了参考帧误差对跟踪稳定性的影响，并展示了 FEMR 对受污染参考动作的修复效果

	In summary, our contributions are:

	\begin{enumerate}
		\item We formulate reference-frame artifacts as a \textit{kinematic-to-dynamic gap} in general motion tracking, and introduce the \textit{robustness budget} to analyze how upstream reference quality governs downstream tracking stability.

		\item We propose FEMR, an online lightweight front-end refiner that keeps the tracker frozen, predicts per-frame task-space corrections gated by confidence scores, and is trained with a tracker-aware executable reward.

		\item Through disturbance recovery experiments, we show that reference-frame noise systematically degrades tracking stability, and that FEMR restores executability without modifying the underlying tracker.
	\end{enumerate}

%===============================================================================

\section{Related Works}

\subsection{General Motion Tracking for Humanoids}

	% Inport 基于强化学习的人形运动跟踪取得了快速进展。本节回顾这一方向：从早期对单一动作的模仿学习，到在多样化数据集上训练的统一全身策略，再到面向通用运动跟踪的规模化尝试

	% PHC 使用大规模动捕数据训练对抗运动先验（AMP），将人体关节运动映射到物理仿真角色。首次以单一策略实现了多样化人体动作的实时全身跟踪

	% ExBody2 将全身控制分解为两部分：上半身跟踪人体动捕的关节角度，下半身通过 RL 独立跟踪根速度指令。双腿不被强制匹配参考轨迹，而是自主决定如何维持平衡

	% OmniH2O 使用蒸馏策略将 VR 头显、动捕服、RGB 相机等多模态输入统一映射到同一控制空间，单一策略同时服务实时遥操作和离线自主执行

	% HumanPlus 构建了全栈模仿管线：WHAM 从单 RGB 相机实时估计人体姿态，GMR 重定向到机器人，Transformer 零样本执行。首次实现了"see-to-do"，但未处理视频姿态估计引入的伪影
	
	% Any2Track 由 AnyTracker 和 AnyAdapter 两级组成：AnyTracker 使用规范动作空间和专家蒸馏训练通用跟踪器；AnyAdapter 通过世界模型从交互历史中提取动力学嵌入，在冻结策略上施加在线自适应，使机器人在飞踹、拉拽、复杂地形等复合扰动下保持稳定

	% KungfuBot2 通过分阶段课程学习逐步扩展策略能力——从静态动作到高动态武术动作逐级训练，最终实现踢腿、旋转跳跃等动作的分钟级连续执行

	% RobotDancing 的策略不预测绝对关节指令，而是输出对参考关节目标的残差修正 Δq。训练时使用分布感知+失败感知的采样策略，确保长尾高动态动作（跳跃、旋转、侧手翻）得到充分覆盖，最终实现多分钟连续舞蹈的零样本 sim-to-real 迁移

	% AMO 在训练中动态调整不同身体部位的跟踪精度权重和物理约束强度，让策略在"精确跟踪"和"物理可行"之间学习自动平衡，而非依赖人工设定的固定奖励权重

	% HOVER 使用单一 Transformer 策略同时支持关节位置、关键点姿态、根速度指令等多种控制模式。不同模式共享运动编码空间，使经验可以跨模式迁移

	% GMT 引入了两项技术：自适应采样根据动作难度自动平衡训练数据分布；Motion MoE 让不同运动区域由不同专家子网络处理。是本文冻结的 tracker

	% UniTracker 采用教师-学生流水线：教师策略在仿真中用特权信息学习隐式运动表示，学生策略通过蒸馏仅凭本体感受复现同等质量的跟踪。隐式运动编码使跟踪器无需逐帧显式参考即可理解当前运动意图

	% SONIC 通过融合多源动捕数据、扩大模型容量、支持多种运动输入接口，探索从规模化跟踪到通用人形控制的技术路线。当前数据规模最大的开源全身跟踪器

	% Export 上述工作证明通用运动跟踪器已不再是动作多样性的瓶颈。但它们共享一个前提：参考动作本身是动力学可行的。当参考来自野外视频、在进入跟踪器前已包含空间伪影时，跟踪器自身的鲁棒性不足以弥补上游信号的质量损失。本文的出发点正在于此——这一前提并不总是成立

	RL-based humanoid motion tracking has advanced rapidly, from imitating individual motions to unified whole-body policies to recent general-purpose tracking systems.
	
	PHC~\cite{phcluo2023} trains an adversarial motion prior (AMP~\cite{amp2021}) on large-scale mocap to map human motion to simulated characters - the first real-time whole-body tracker. ExBody2~\cite{exbodyji2024} decouples upper-body pose tracking from lower-body RL velocity control, letting legs decide how to balance independently. OmniH2O~\cite{omnih2ohe2024} distills multimodal inputs (VR, mocap, RGB) into a shared control space for both teleoperation and autonomous execution. HumanPlus~\cite{humanplus2024} builds a full-stack pipeline - WHAM~\cite{whamshin2024}, GMR~\cite{gmr2025}, Shadowing Transformer - for zero-shot see-to-do, but does not address artifacts from video pose estimation. Any2Track~\cite{zhang2025track} pairs a generalist tracker (AnyTracker) with a history-conditioned adaptation module (AnyAdapter) for robustness under combined physical disturbances. KungfuBot2~\cite{han2025kungfubot2} extends policy capability from static to highly dynamic motions via staged curriculum learning. RobotDancing~\cite{sun2025robotdancing} predicts residual joint corrections $\Delta q$ on reference targets, using distribution- and failure-aware sampling for long-tail dynamic motions. AMO~\cite{li2025amo} dynamically adjusts per-body-part tracking weights during training to balance precision and feasibility. HOVER~\cite{hover2025} supports joint, keypoint, and velocity control modes within a single Transformer via a shared motion encoding space. GMT~\cite{gmtchen2025} combines adaptive motion sampling with a Motion MoE architecture. It is the frozen tracker in this work. UniTracker~\cite{uniyin2025} distills privileged teacher information into a proprioception-only student via latent motion encoding. SONIC~\cite{sonic2025} scales up tracking via multi-source mocap fusion and increased model capacity, the largest open-source tracker to date.
	
	These works show that general trackers are no longer the bottleneck for motion diversity. But they share a premise: the reference is dynamically executable. When references come from in-the-wild videos with spatial artifacts already present, the tracker's robustness is insufficient to compensate. This paper starts from the observation that this premise does not always hold.

\subsection{From In-the-Wild Videos to Reference Motions}

	% Inport 通用跟踪器的能力提升使参考动作的来源成为新的关注点。从单目视频恢复人体的三维姿态和全局轨迹，是获取可扩展参考动作的关键途径。本节回顾这一方向：从单帧人体网格恢复，到在世界坐标系中估计完整时序运动

	% HMR 使用对抗学习从单目 RGB 图像直接回归 SMPL 人体模型的姿态和体型参数，结合 2D 关键点重投影损失和 3D 对抗先验，无需成对的 2D-3D 监督。是单目人体重建领域的奠基工作。其输出的深度歧义和姿态误差会沿整个管线向下传递——这是伪影的最上游来源
	
	% GLAMR 在动态相机下联合优化相机运动、人体姿态和全局轨迹，使用运动先验填补遮挡期间的人体轨迹，避免遮挡导致身份切换或轨迹中断
	
	% TRACE 在一个端到端网络中同时恢复和跟踪动态相机视频中的多人 3D 世界坐标轨迹。它用 World Motion Map 将人体运动从相机坐标映射到全局坐标，首次统一解决了多人跟踪和全局轨迹恢复，无需分两阶段处理
	
	% WHAM 融合 AMASS 运动先验和 SLAM 估计的相机角速度，通过运动编码器将 2D 关键点序列提升为 3D 表征，再与视频图像特征融合，将人体从相机坐标恢复至全局世界坐标。首次让视频方法在逐帧精度上超越所有单帧方法。其全局轨迹精度受限于 SLAM——在动态人体遮挡场景下相机运动估计不可靠

	% GVHMR 用重力-视角坐标替代绝对世界坐标：以重力方向为固定参考轴，仅估计人体相对重力的旋转和平移。在倾斜相机、上下坡等场景下显著减少了根节点朝向漂移。但重力方向本身需要从视频中估计，在快速运动和无地平线参照时仍有不可消除的误差
	
	% TRAM 分两步恢复人体的全局轨迹和运动：先鲁棒化 SLAM 从静态背景中恢复度量尺度的相机轨迹，再通过视频 Transformer（VIMO）回归人体运动。它对动态人体区域做掩码处理，使 SLAM 仅依赖静态背景。但在背景纹理稀疏或高度动态场景下 SLAM 仍可能发散
	
	% PHMR 将文本或 2D 视觉提示作为条件信号注入重建网络，用 CLIP 嵌入空间对齐文本描述和人体姿态特征，使重建结果不仅准确而且符合语义约束

	% Export 上述方法产出的三维人体运动在姿态空间中通常是合理的。但单目重建的固有歧义使得全局轨迹中不可避免地残留空间误差：根节点高度漂移、绝对朝向偏移、水平位移累积。这些误差在视觉上不易察觉，却足以改变跟踪器依赖的物理条件。现有修正手段依赖离线优化和完整序列——它需要全部未来帧和可观的求解时间。当目标是零样本即时执行时，离线修正天然不适用。当前管线在运动学重建与跟踪器之间，缺少一个在线、逐帧、不依赖完整序列的前端修正环节

	As trackers improve, attention shifts upstream. Recovering 3D human pose and global trajectory from monocular video is key to scalable reference acquisition - we review the progression from single-frame mesh recovery to full temporal motion in world coordinates. HMR~\cite{hmr2018} regresses SMPL parameters from a single RGB image via adversarial learning. Its depth ambiguity is the most upstream source of artifacts in the pipeline. GLAMR~\cite{yuan2022glamr} jointly optimizes camera motion, body pose, and global trajectory under dynamic cameras, using a motion prior to fill occlusions. TRACE~\cite{sun2023trace} recovers multi-person 3D world trajectories in a single end-to-end network via its World Motion Map. WHAM~\cite{whamshin2024} fuses an AMASS motion prior with SLAM-estimated camera velocity to recover world-grounded humans; its accuracy is bounded by SLAM reliability under occlusion. GVHMR~\cite{shen2024world} anchors the body to a gravity direction axis to reduce orientation drift under tilted cameras, though gravity estimation itself remains imperfect. TRAM~\cite{tramwang2024} robustifies SLAM on static backgrounds, then regresses body motion with a video Transformer; SLAM diverges in sparse or highly dynamic scenes. PHMR~\cite{wang2025prompthmr} injects text or visual prompts into reconstruction via CLIP embeddings for semantically consistent output.

	These methods produce motions that are plausible in pose space but not physically validated. Monocular ambiguity leaves inevitable spatial errors - root drift, orientation offset, horizontal displacement - that alter the physical conditions trackers depend on. No online, per-frame correction stage currently exists between kinematic reconstruction and the tracker.

	% PHC 将 SMPL 人体姿态通过关键点匹配映射到人形机器人关节空间，是 RL-based tracking 中早期常用的重定向方案
	% PHC~\cite{phcluo2023} maps SMPL human poses to humanoid joint space via keypoint matching - an earlier retargeting scheme in RL-based tracking.

	% GMR 使用两阶段 IK 优化实现人体到机器人的运动重定向：先优化根节点朝向消除初始位姿偏差，再做全身微调。其关键创新是非均匀局部缩放——对不同身体部位使用独立的缩放因子
	% GMR~\cite{gmr2025} uses two-stage IK optimization with non-uniform per-body-part scaling for higher-quality retargeting.

	% OmniRetarget 在机器人、物体、地面之间构建交互网格，通过拉普拉斯形变能量在重定向过程中保持空间交互关系不变，同时施加硬物理约束（关节限位、自碰撞、足底接触）。它能保证输出轨迹物理自洽，但无法修正输入人体姿态中已存在的空间误差
	% OmniRetarget~\cite{yang2025omniretarget} preserves robot-object-ground interaction via an interaction mesh with hard physics constraints. It guarantees self-consistent output but cannot correct errors already in the input pose.

	% ProtoMotions 从动捕数据提取有限个原型姿态作为动作基，用对抗运动先验将真实人体动作投影到原型组合上。在重定向管线中作为可选后处理步骤，对违反物理约束的关节轨迹做平滑修正
	% ProtoMotions~\cite{ProtoMotions} projects motions onto a learned prototype basis as post-processing for physics violations.

\subsection{Residual Learning and Sim-to-Real Adaptation}

	% Inport 残差学习提供了一种无需重新训练主干网络的修正范式：冻结已有的基座控制器，在其上训练轻量的加性修正模块。本节回顾这一范式从机器人操作到足式运动、再到人形全身控制中的发展

	% Residual RL 率先提出残差强化学习范式：冻结手写反馈控制器，在其上训练 RL 策略输出加性残差动作。基座策略负责已知动力学，残差策略仅学习接触密集和不确定性区域
	
	% RMA 将策略分解为基座和自适应两部分：基座在多样化仿真环境（不同摩擦、地形、载荷）中训练后冻结；自适应模块通过监督学习从本体感受历史实时推断环境隐变量并调制基座行为。首次在真实足式机器人上展示了在线自适应多地形能力
	
	% ASAP 先训练 motion-tracking 策略，再用真实数据训练 delta action model 补偿仿真动力学失配，最后将 delta model 蒸馏回基座策略。delta model 被冻结后作为仿真环境的动力学修正层，基座策略在该增强仿真中微调，部署时不再需要 delta model。实现了跳跃、后空翻等敏捷动作的零样本迁移

	% ResMimic 分两阶段训练：先预训练通用跟踪策略作为运动先验，再冻结该先验、针对特定操作任务训练残差策略。残差仅学习操作相关的关节修正，不干扰已学好的运动技能。残差施加在动作空间

	% MOSAIC 冻结通用跟踪策略，用少量遥操作数据训练轻量残差模块，通过双教师蒸馏将残差修正注入跟踪器同时防止灾难性遗忘。与 FEMR 架构最接近——冻结 tracker + 轻量修正——但修正施加在动作空间
	
	% OmniXtreme 使用 flow matching 从大规模动捕数据生成运动先验分布，再通过残差 RL 后训练弥合仿真-真实差距。用生成模型的结构化运动先验替代传统轨迹跟踪参考信号，使策略在行为空间而非轨迹空间中探索。残差仍作用于 tracker 输出端

	% Export 残差学习的有效性已在多种机器人任务中得到验证。但在人形运动控制中，已有残差方法均将修正施加在动作空间——它们改变的是跟踪器的输出。当问题源头在上游参考信号而非跟踪器自身时，在参考空间施加修正更为直接。在参考层做在线、逐帧的前端残差修正——这一方向尚未被探索

	Residual learning avoids retraining the backbone: freeze a base controller and train a lightweight additive correction module on top. We trace this paradigm from manipulation to legged locomotion to humanoid whole-body control.
	
	Residual~RL~\cite{johannink2019residual} first formalized the paradigm: a frozen base controller plus a learned additive residual. FEMR shares this architecture but applies residuals in reference rather than action space. RMA~\cite{kumar2021rma} splits into a base policy (frozen after diverse training) and an adaptation module that infers environment latents from proprioceptive history. ASAP~\cite{asap2025} trains a delta action model on real data to correct sim dynamics mismatch, then distills it back into the base for deployment. ResMimic~\cite{resmimic2025} pretrains a generalist tracking prior, then freezes it and trains task-specific residuals in action space. MOSAIC~\cite{mosaic2026} freezes a tracking policy and trains a residual module on teleoperation data via dual-teacher distillation. OmniXtreme~\cite{omnixtreme2026} uses flow matching to generate motion priors, then applies residual RL post-training. The residual still acts on the tracker output.

	Residual learning has been validated across diverse tasks. In humanoid motion control, however, existing residual methods all correct in action space - they modify the tracker's output. When the problem is upstream in the reference signal, correcting in reference space is more direct. Online, per-frame, front-end residual correction at the reference level remains unexplored.

%===============================================================================

\section{Method}





%===============================================================================

\section{Citations}
\label{sec:citations}

	Citations can be made using either \textbackslash citep\{\} or \textbackslash citet\{\}, depending from the appropriateness. To avoid the citation moving to the next line, it is often a good practice to replace the space before with a tilde (\~{}) character.
	Example 1: ``CoRL is the best conference ever~\citep{Gauss1857}.''
	Example 2: ``\citet{Lagrange1788} proved, both theoretically and numerically, that CoRL is the best conference ever.''
	
%===============================================================================

\section{Experimental Results}
\label{sec:result}

	

%===============================================================================

\section{Conclusion}
\label{sec:conclusion}

	

%===============================================================================

\clearpage
% The acknowledgments are automatically included only in the final and preprint versions of the paper.
\acknowledgments{If a paper is accepted, the final camera-ready version will (and probably should) include acknowledgments. All acknowledgments go at the end of the paper, including thanks to reviewers who gave useful comments, to colleagues who contributed to the ideas, and to funding agencies and corporate sponsors that provided financial support.}

%===============================================================================

% no \bibliographystyle is required, since the corl style is automatically used.
\bibliography{example}  % .bib

\end{document}
